\begin{afterword}
\summaryhead

A campaign manager, newly hired by a mayoral candidate, asks how to build an impeccable rational argument that the candidate is best. The post answers that it cannot be done, even if every fact cited is true and relevant, because the conclusion was fixed in advance.

The example is a newspaper questionnaire on which the candidate does well on 15 questions out of 16. The one exception, question 11, asks whether he is now or has ever been a supervillain. Publishing the questionnaire without question 11 would force voters to infer hidden evidence, and the manager crossed the line as soon as publication depended on whether the answers helped.

The post then notes that an argument can be invalid even when its premises and conclusion are true, and calls an argument whose stated reasons are not its real reasons ``hypocritical.'' The only honest method is to choose a candidate from the evidence before anyone hires you, and then print your checklist as campaign literature.

\responsehead

The questionnaire is a good example. Every answer the campaign would publish is true, and the result would still mislead, because a voter who knows how the campaign chooses what to print must suspect what is missing. This applies ``What Evidence Filtered Evidence?'' from three days before.

The thesis is stated more broadly than the example supports. ``It can't be done'' is true in the sense the post gives ``rational argument'' in its last paragraph, an argument whose conclusion was written after its premises. That sense is fixed by the post's own definition. In the sense the manager means, an argument from true and relevant premises, the claim holds for arguments that pick among evidence, where what is left out changes what the rest shows. It does not hold for deduction. A mathematician often knows the theorem before finding the proof, and the proof is valid or not whatever the order of writing. When a reader raised this in the comments, the author granted that an identical text ``might be just as good'' whoever wrote it, and answered that ``there's a reason why it wasn't.''

The syllogism illustrates a different fault from the one the post is about. Its premises do not imply its conclusion, so its stated reasons are no reasons at all. The campaign's facts are, on the post's own terms, relevant evidence, and the problem is what was left out. The post uses the first to introduce a gap between stated and real reasons, but only the second is the problem of the questionnaire.

The closing rule, that an honest arguer can write above the conclusion only what actually decided it, is about the honesty of the arguer. A reader should not take it to mean that facts cited by an interested party are worthless. The post does not say that, and in a comment the same day the author made the point directly: ``Valid evidence is valid, whatever the motives of the one who cites it.''

In short: a good example of how true facts, selectively published, mislead, together with a thesis that is true by the post's definition and too broad in the ordinary sense, and to a syllogism that shows a different error.
\end{afterword}
